

\chapter{REFERENCIAL TEÓRICO}


% \section{INDÚSTRIA DE JOGOS DIGITAIS}

% A indústria de jogos digitais se tornou, nas últimas décadas, um dos setores mais expressivos da economia do entretenimento mundial, superando em receita outros segmentos tradicionais, como o cinema e a música, tendo movimentado bilhões de dólares, e sua trajetória impulsionada pela popularização de plataformas de distribuição digital e pela diversificação dos modelos de monetização (PASHKOV, 2021).

% Esse crescimento não se limita apenas ao aspecto financeiro, refletindo também em mudanças estruturais na forma como os jogos são produzidos e consumidos globalmente. No contexto brasileiro, de acordo com Pesquisa da Indústria Brasileira de Games 2022 a indústria de jogos digitais, embora ainda em processo de amadurecimento em comparação a mercados internacionais consolidados, apresenta um ecossistema crescente e diversificado, impulsionado por políticas de fomento, editais de financiamento e pela formação de polos regionais de desenvolvimento  \cite{fortim2022pesquisa}.

% Diante desse cenário de expansão global e nacional, a colaboração entre desenvolvedores e comunidade de jogadores torna-se um elemento cada vez mais relevante no ciclo de desenvolvimento.


\section{DESENVOLVIMENTO COLABORATIVO}

Segundo \citeonline{uchoa2018} o crowdsourcing no desenvolvimento de jogos vai além de apenas coletar opiniões, ele se configura como um modelo de produção no qual a comunidade participa ativamente da construção do produto. Nesse contexto, Linåker, Bjarnason e Fagerholm (2025) destacam que a validação contínua de ideias junto ao público tornou-se uma estratégia importante para reduzir riscos e orientar decisões no projeto, especialmente em contextos nos quais desenvolvedores operam com prazos curtos e recursos limitados.

Complementando essa visão, \citeonline{thominet2018} discute o conceito de Desenvolvimento Aberto (Open Development) como uma metodologia que transforma a comunicação como elemento central no ciclo de produção, e não apenas algo secundário. De acordo com o autor, o compromisso com a transparência e o acesso contínuo da comunidade ao desenvolvimento não apenas melhora o produto, mas também constrói uma base de usuários engajada.

\citeonline{moro2022} ampliam essa discussão ao demonstrar que a qualidade e volume do feedback obtido via comunidades digitais muitas vezes supera a de testes em ambientes controlados, pois reflete experiências reais de uso em contextos variados. \citeonline{tong2022}, por sua vez, alerta que a dispersão do feedback entre diferentes plataformas representa um obstáculo, visto que cada ambiente produz informações de maneiras distintas e com diferentes graus de profundidade, o que exige dos desenvolvedores uma categorização estratégica das contribuições recebidas.


\section{DESIGN DE INTERFACE E USABILIDADE}
